\begin{lemma}[Continuity of a metric $1$-current with respect to pointwise convergence of $f$]\label{CurrentContinuousF}
  Let $E$ be a metric space equipped with a compatible Borel $\sigma$-algebra, and $T \in
  \mathcal{M}_1(E)$ a metric $1$-current (\noderef{MetricCurrent1}). Fix $\pi \colon E \to
  \mathbb{R}$ Lipschitz with constant $\mathrm{Lip}(\pi)$, and $f \colon E \to \mathbb{R}$ bounded
  ($\norm{f}_\infty \le b_F$) and Lipschitz (constant $L_F$). Let $(f_n)_{n \in \mathbb{N}}$ be a
  sequence of bounded Lipschitz functions, bounded ($\norm{f_n}_\infty \le b_n$) and Lipschitz
  (constant $L_n$), with the bounds $(b_n)$ uniformly bounded by some $b_0$, such that
  $f_n(x) \to f(x)$ as $n\to\infty$ for every $x \in E$. Then
  \[
    T(f_n,\pi) \to T(f,\pi) \quad \text{as } n \to \infty
    ,
  \]
  where $T(f_n,\pi)$ and $T(f,\pi)$ denote $T$ applied to the metric $1$-form (\noderef{Form1})
  with the indicated $f$- and $\pi$-components (and the stated witnessing bound/Lipschitz data).

  \paragraph{Provenance and why this is a new node.} Paper.tex line 1393, in the proof of
  Theorem~4.4, states: ``we recall that by definition, a metric $1$-current $T$ is continuous with
  respect to the pointwise convergence [of $\pi$, eq.~\eqref{pi_converge}], while the continuity
  with respect to [the analogous pointwise convergence of $f$, eq.~\eqref{f_converge}] follows from
  [the mass-measure inequality, eq.~\eqref{massmeasure}] and Lebesgue's dominated convergence
  theorem.'' \noderef{MetricCurrent1}'s own docstring records the same fact and, correspondingly,
  \emph{excludes} this statement from that structure's fields: continuity in $\pi$
  (\texttt{continuous\_pi}) is axiomatized directly, but continuity in $f$ is deliberately left
  underivable from the structure fields alone, to be derived at whichever node needs it, precisely
  because the paper itself derives it (rather than positing it) from finite mass together with
  dominated convergence. This lemma is exactly that derivation, stated once as its own node so
  that \noderef{DiagonalizationGeneral} (which needs it as one of the two halves of
  eq.~\eqref{TfpiTfkpij}, the other half being \noderef{MetricCurrent1}'s \texttt{continuous\_pi}
  directly) can cite a single clean statement rather than repeating a genuine Lebesgue
  dominated-convergence argument (using \texttt{finiteMass}, additivity in the $f$-slot, and the
  admissible-measure domination bound) inline. This is not a Substantiveness Clause-2 wrapper: it
  is not a repackaging of any single existing node's content into a shorter proof of the same fact,
  since no existing tablet node states or proves it, and its own proof requires a genuine
  application of the dominated convergence theorem, not a one- or few-line consequence of a single
  cited node.
\end{lemma}
\begin{proof}
  \emph{Step 1: exhibit an admissible measure.} By \noderef{MetricCurrent1}'s \texttt{finiteMass}
  field, $\mathrm{massOfFunctional}(T) \ne \infty$ (\noderef{massOfFunctional}). By
  \noderef{massOfFunctional}'s definition as an infimum over the subtype of measures admissible for
  $T$ (\noderef{IsAdmissible}), if that subtype were empty the infimum over it would be $\infty$
  (the infimum of the empty set in $[0,\infty]$, by convention), contradicting
  \texttt{finiteMass}. Hence the subtype is nonempty: there exists a measure $\mu$ on $E$ that is
  admissible for $T$, i.e.\ (\noderef{IsAdmissible}) $\mu$ is a finite Borel measure and
  \[
    |T(\omega)| \le \mathrm{Lip}(\omega.\pi) \int_E |\omega.f| \, d\mu
    \qquad \text{for every } \omega \in \mathcal{D}^1(E)
    . \tag{$\ast$}
  \]
  Fix such a $\mu$ for the rest of the proof.

  \emph{Step 2: split off the difference form.} For each $n$, let $D_n \in \mathcal{D}^1(E)$ be the
  metric $1$-form with $f$-part $x \mapsto f_n(x) - f(x)$ (bound $b_n + b_F$, by the triangle
  inequality $|f_n(x)-f(x)| \le |f_n(x)|+|f(x)| \le b_n+b_F$; Lipschitz constant $L_n+L_F$, since a
  difference of an $L_n$-Lipschitz and an $L_F$-Lipschitz function is $(L_n+L_F)$-Lipschitz) and
  $\pi$-part $\pi$ unchanged (Lipschitz constant $\mathrm{Lip}(\pi)$). Since $f_n(x) = (f_n(x)-
  f(x)) + f(x)$ pointwise and $D_n.\pi = \pi = (f,\pi).\pi$ and $(f_n,\pi).\pi = \pi$ all coincide,
  \noderef{MetricCurrent1}'s \texttt{add\_f} field (additivity in the $f$-slot, applied with
  $\omega_{12} := (f_n,\pi)$, $\omega_1 := D_n$, $\omega_2 := (f,\pi)$) gives
  \[
    T(f_n,\pi) = T(D_n) + T(f,\pi)
    \qquad \text{for every } n
    . \tag{$\ast\ast$}
  \]
  It therefore suffices to show $T(D_n) \to 0$ as $n \to \infty$.

  \emph{Step 3: the $L^1(\mu)$-norm of $D_n.f$ vanishes, by dominated convergence.} The function
  $x \mapsto |D_n.f(x)| = |f_n(x)-f(x)|$ is measurable for each $n$ (it is continuous, being built
  from the Lipschitz, hence continuous, functions $f_n,f$). It is dominated uniformly in $n$ by the
  constant $b_0+b_F$: for every $x$, $|f_n(x)-f(x)| \le |f_n(x)|+|f(x)| \le b_n+b_F \le b_0+b_F$
  (using $b_n \le b_0$ for every $n$, the hypothesis). This constant is $\mu$-integrable since $\mu$
  is a finite measure (Step 1). And for every fixed $x \in E$, $|f_n(x)-f(x)| \to |f(x)-f(x)| = 0$
  as $n \to \infty$, since $f_n(x) \to f(x)$ (the hypothesis) and subtraction and the absolute value
  are continuous. By Lebesgue's dominated convergence theorem (in its $[0,\infty]$-valued lintegral
  form, \texttt{MeasureTheory.tendsto\_lintegral\_of\_dominated\_convergence}, \texttt{Preamble}),
  applied with the dominating function the constant $b_0+b_F$ and limit function identically $0$,
  \[
    \int_E |D_n.f| \, d\mu \longrightarrow 0 \qquad \text{as } n \to \infty
    . \tag{$\dagger$}
  \]

  \emph{Step 4: squeeze $T(D_n)$ to $0$.} Multiplying $(\dagger)$ by the (finite) constant
  $\mathrm{Lip}(\pi)$ preserves the limit:
  \[
    \mathrm{Lip}(\pi) \int_E |D_n.f| \, d\mu \longrightarrow 0
    \qquad \text{as } n \to \infty
    . \tag{$\ddagger$}
  \]
  By the domination bound $(\ast)$ applied to $\omega := D_n$ (whose $\pi$-part is $\pi$, so
  $\mathrm{Lip}(D_n.\pi) = \mathrm{Lip}(\pi)$),
  \[
    0 \le |T(D_n)| \le \mathrm{Lip}(\pi) \int_E |D_n.f| \, d\mu
    \qquad \text{for every } n
    .
  \]
  Since the left bound is the constant sequence $0$ (which trivially tends to $0$) and the right
  bound tends to $0$ by $(\ddagger)$, the squeeze theorem gives $|T(D_n)| \to 0$, and hence
  $T(D_n) \to 0$, as $n \to \infty$.

  \emph{Step 5: conclude.} Combining Step 4 with $(\ast\ast)$ and continuity of addition,
  \[
    T(f_n,\pi) = T(D_n) + T(f,\pi) \longrightarrow 0 + T(f,\pi) = T(f,\pi)
    \qquad \text{as } n \to \infty
    ,
  \]
  which is the claimed limit.
\end{proof}
